% part: intuitionistic-logic
% chapter: introduction
% section: bhk-interpretation

\documentclass[../../../include/open-logic-section]{subfiles}

\begin{document}

\olfileid{int}{int}{bhk}

\olsection{ब्रौवर--हीटिंग--कोल्मोगोरोव अर्थनिर्धारण}

\begin{editorial}
  BHK अर्थनिर्धारण वापरून अंतःप्रज्ञावादी विधानांची
  वैधता सिद्ध करणाऱ्या सिद्धता गोंधळात टाकतात;
  त्यांचे अधिक चांगले स्पष्टीकरण द्यायला हवे.
\end{editorial}

अंतःप्रज्ञावादी संयोजकांचे एक अनौपचारिक रचनाशील
अर्थनिर्धारण आहे; त्याला बहुधा ब्रौवर--हीटिंग--कोल्मोगोरोव
अर्थनिर्धारण म्हणतात. त्यात ``रचना'' ही संकल्पना
वापरली जाते; तिला रचनाशील सिद्धता समजू शकता.
(औपचारिक !!{derivation} व्यवस्थेतील
!!a{derivation} शी गल्लत होऊ नये म्हणून BHK
अर्थनिर्धारणात ``सिद्धता'' हा शब्द वापरत नाही.)
या सहज समजणाऱ्या संकल्पनेच्या आधारे BHK अर्थनिर्धारण
अंतःप्रज्ञावादी संयोजकांचे अर्थ स्पष्ट करते.

\begin{enumerate}
\item अण्विक विधानाची रचना कशाला म्हणायचे हे आपल्याला
  माहीत आहे असे गृहीत धरू.
\item $!A_1 \land !A_2$ ची रचना म्हणजे $\tuple{M_1, M_2}$
  ही जोडी; येथे $M_1$ ही~$!A_1$ ची आणि $M_2$ ही~$!A_2$
  ची रचना आहे.
\item $!A_1 \lor !A_2$ ची रचना म्हणजे $\tuple{s, M}$
  ही जोडी; यात $s$ हे~$1$ आणि $M$ ही~$!A_1$ ची रचना
  असते, किंवा $s$ हे~$2$ आणि $M$ ही~$!A_2$ ची रचना
  असते.
\item $!A \lif !B$ ची रचना म्हणजे~$!A$ ची रचना
  घेऊन तिला~$!B$ च्या रचनेत रूपांतरित करणारे फलन.
\item $\lfalse$ (असंगती) ची कोणतीही रचना नसते.
\item $\lnot !A$ हे $!A \lif \lfalse$ चे समानार्थी रूप
  म्हणून परिभाषित केले आहे. म्हणजे $\lnot !A$ ची रचना
  ही~$!A$ ची रचना घेऊन तिला~$\lfalse$ च्या रचनेत
  रूपांतरित करणारे फलन असते.
\end{enumerate}

\begin{ex}
उदाहरणार्थ $\lnot \lfalse$ घ्या. त्याची रचना हे
असे फलन असते जे~$\lfalse$ ची कोणतीही रचना इनपुट
म्हणून घेऊन~$\lfalse$ ची रचना आउटपुट म्हणून देते.
ओळख फलन~$\Id{}$ ही अशीच एक रचना आहे:
$\lfalse$ ची रचना~$M$ दिली, तर $\Id{}(M) = M$
ही~$\lfalse$ ची रचना देते.
\end{ex}

साधारणपणे $\lnot !A$ चा अर्थ ``$!A$ ची रचना अशक्य आहे'' असा होतो.

\begin{ex}
कोणत्याही विधान~$!A$ साठी $!A \lif \lnot \lnot !A$
सिद्ध करू. हेच सूत्र $!A \lif ((!A \lif \lfalse) \lif
\lfalse)$ आहे. याची रचना असे फलन~$f$ असली पाहिजे
की $!A$ ची रचना~$M$ दिल्यावर ते $(!A \lif \lfalse)
\lif \lfalse$ ची रचना~$f(M)$ देते. हे फलन~$f$
$(!A \lif \lfalse) \lif \lfalse$ ची रचना कशी
तयार करते ते पाहू: $!A \lif \lfalse$
ची रचना~$h$ इनपुट म्हणून घेऊन~$\lfalse$ ची रचना
आउटपुट म्हणून देणारे फलन $g$ परिभाषित करायचे आहे.
त्यासाठी $g$ ची व्याख्या अशी करा: इनपुट~$h$ हे
आधी मिळालेल्या~$!A$ च्या
रचनेवर~$M$ लावा. $h$ चे आउटपुट~$h(M)$ ही
$\lfalse$ ची रचना असल्याने, $M$ ही~$!A$ ची रचना
असेल तेव्हा $f(M)(h) = h(M)$ ही~$\lfalse$ ची रचना
असते.
\end{ex}

\begin{ex}
आता $\lnot(!A \land \lnot !A)$, म्हणजेच $(!A
\land (!A \lif \lfalse)) \lif \lfalse$, याची रचना
देऊ. ती असे फलन~$f$ आहे की $!A \land (!A \lif
\lfalse)$ ची रचना~$M$ इनपुट म्हणून दिल्यास ते
$\lfalse$ ची रचना देते. संयोग $!B_1 \land !B_2$
ची रचना म्हणजे $\tuple{N_1, N_2}$ ही जोडी; येथे
$N_1$ ही~$!B_1$ ची आणि $N_2$ ही~$!B_2$ ची
रचना आहे. $!B_1 \land !B_2$ च्या रचनेतून
अनुक्रमे $!B_1$ आणि $!B_2$ यांच्या रचना मिळवणारी
फलने $p_1$ आणि $p_2$ अशी परिभाषित करू:
\begin{align*}
  p_1(\tuple{N_1, N_2}) &= N_1\\
  p_2(\tuple{N_1, N_2}) &= N_2
\end{align*}
फलन~$f$ पुढीलप्रमाणे काम करते. प्रथम ते आपल्या
इनपुट~$M$ ला $p_1$ लावते; त्यामुळे~$!A$ ची रचना
मिळते. नंतर~$M$ ला $p_2$ लावते; त्यामुळे
$!A \lif \lfalse$ ची रचना मिळते. ही रचना म्हणजे
फलन~$p_2(M)$; त्याला~$!A$ ची रचना इनपुट म्हणून
दिल्यास ते~$\lfalse$ ची रचना देते. दुसऱ्या शब्दांत,
$p_2(M)$ हे $p_1(M)$ ला लावल्यास~$\lfalse$ ची
रचना मिळते. म्हणून $f(M) = p_2(M)(p_1(M))$ अशी
व्याख्या करता येते.
\end{ex}

\begin{ex}
आता $((!A \land !B) \lif !C) \lif (!A \lif
(!B \lif !C))$ ची रचना देऊ. म्हणजे, $(!A \land
!B) \lif !C$ ची रचना~$g$ घेऊन $(!A \lif (!B \lif
!C))$ ची रचना देणारे फलन $f$ हवे. रचना~$g$
ही स्वतः फलन आहे: $!A \land !B$ च्या रचनांपासून
$!C$ च्या रचनांकडे जाणारे. आणि आउटपुट~$f(g)$
हे फलन~$h_g$ आहे: $!A$ च्या रचनांपासून, $!B$
च्या रचनांना~$!C$ च्या रचनांत नेणाऱ्या फलनांकडे जाणारे.

ठीक आहे, हे गोंधळात टाकणारे आहे. इनपुट~$g$ साठी
$f$ चे आउटपुट ठरणारे फलन $h_g$ तयार करायचे आहे.
$h_g$ चा इनपुट ही~$!A$ ची रचना~$M$ आहे.
$h_g(M)$ चे आउटपुट हे~$!B$ च्या रचनांपैकी~$N$
पासून~$!C$ च्या रचनांकडे जाणारे फलन~$k_M$ असले
पाहिजे. $k_{g,M}(N) = g(\tuple{M, N})$ असे ठेवू.
$\tuple{M, N}$ ही~$!A \land !B$ ची रचना आहे हे
आठवा. म्हणून $k_{g,M}$ ही $!B \lif !C$ ची रचना
आहे: ती~$!B$ च्या रचना~$N$ यांना~$!C$ च्या
रचनांत नेते. आता $h_g(M) = k_{g,M}$ असे ठेवू.
हे फलन~$!A$ च्या रचना~$M$ यांना $!B
\lif !C$ च्या रचना~$k_{g,M}$ यांत नेते. पुढे
$f(g) = h_g$ असे ठेवू. हे फलन $(!A \land !B)
\lif !C$ च्या रचना~$g$ यांना $!A \lif (!B \lif !C)$
च्या रचनांत नेते. हुश्श!{}
\end{ex}

$!A \lor \lnot !A$ या विधानाला विमध्य सिद्धांत
म्हणतात. काही विशिष्ट~$!A$ साठी (उदाहरणार्थ,
$\lfalse \lor \lnot \lfalse$) तो सिद्ध करता येतो;
पण सर्वसाधारणपणे नाही. कारण अंतःप्रज्ञावादी
विकल्पयोगासाठी दोन विकल्पांपैकी एकाची रचना लागते;
पण काही विधाने अशी आहेत की सध्या ती सिद्धही
करता येत नाहीत किंवा खोडताही येत नाहीत
(उदाहरणार्थ, गोल्डबाखचे अनुमान). मात्र विमध्य
सिद्धांत खोडताही येत नाही: म्हणजे $\lnot \lnot (!A
\lor \lnot !A)$ सत्य असते.

\begin{ex}
$\lnot \lnot (!A \lor \lnot !A)$ सिद्ध करण्यासाठी
$\lnot(!A \lor \lnot !A)$, म्हणजे $(!A
\lor (!A \lif \lfalse)) \lif \lfalse$, याची रचना
घेऊन तिला~$\lfalse$ च्या रचनेत रूपांतरित करणारे
फलन~$f$ हवे. दुसऱ्या शब्दांत, $g$ ही $\lnot(!A
\lor \lnot !A)$ ची रचना असेल, तर $f(g)$ ही
$\lfalse$ ची रचना ठरेल, असे फलन $f$ हवे.

$g$ ही $\lnot(!A \lor \lnot !A)$ ची रचना आहे असे
समजा; म्हणजे ते $!A \lor \lnot !A$ ची रचना घेऊन
तिला~$\lfalse$ च्या रचनेत नेणारे फलन आहे.
$!A \lor \lnot !A$ ची रचना म्हणजे $\tuple{s, M}$
ही जोडी: यात एकतर $s = 1$ आणि $M$ ही~$!A$ ची
रचना असते, किंवा $s = 2$ आणि $M$ ही $\lnot !A$
ची रचना असते. $h_1$ हे~$!A$ ची रचना~$M_1$ घेऊन
$!A \lor \lnot !A$ ची रचना देणारे फलन असू द्या:
ते $M_1$ ला $\tuple{1, M_1}$ कडे नेते. तसेच $h_2$
हे $\lnot !A$ ची रचना~$M_2$ घेऊन
$!A \lor \lnot !A$ ची रचना देणारे फलन असू द्या:
ते $M_2$ ला $\tuple{2, M_2}$ कडे नेते.

$k$ हे $\comp{h_1}{g}$ असू द्या. ते~$!A$ ची रचना
दिल्यावर~$\lfalse$ ची रचना देणारे फलन आहे; म्हणजे
ते~$!A \lif \lfalse$ किंवा $\lnot !A$ ची रचना
आहे. आता $l$ हे $\comp{h_2}{g}$ असू द्या. हे
फलन $\lnot !A$ ची रचना दिल्यावर~$\lfalse$ ची
रचना देते. $k$ ही $\lnot !A$ ची रचना असल्याने,
$l(k)$ ही~$\lfalse$ ची रचना आहे.

अशा प्रकारे $\lnot(!A \lor \lnot !A)$ ची
रचना~$g$ घेऊन तिला~$\lfalse$ च्या रचनेत कसे
रूपांतरित करायचे ते सांगितले. म्हणजे
$\lnot(!A \lor \lnot !A)$ ची रचना~$g$
घेऊन~$\lfalse$ ची रचना~$l(k)$ देणारे फलन $f$
ही $\lnot\lnot(!A \lor \lnot !A)$ ची रचना आहे.
\end{ex}

पाहिल्याप्रमाणे, !!{formula}s ची अंतःप्रज्ञावादी
वैधता दाखवण्यासाठी BHK अर्थनिर्धारण वापरणे लवकरच
अवजड आणि गोंधळात टाकणारे होते. सुदैवाने,
अंतःप्रज्ञावादी तर्कशास्त्रासाठी अधिक चांगल्या
!!{derivation} व्यवस्था आणि अधिक नेमकी
चिन्हार्थमीमांसक अर्थलक्षणे उपलब्ध आहेत.
\end{document}
